\documentclass[11pt]{article}
\usepackage[T1]{fontenc}
\usepackage{lmodern,amsmath,amssymb,amsthm,mathtools,booktabs}
\usepackage[margin=1in]{geometry}
\usepackage[colorlinks=true,linkcolor=blue!45!black,urlcolor=blue!45!black,citecolor=blue!45!black]{hyperref}
\usepackage{xcolor}
\usepackage{microtype}
\setlength{\emergencystretch}{2em}
\hypersetup{pdftitle={Sharp support constants for positive filters of Riemann and Dirichlet beta zeros},pdfauthor={Research note for ProveIt}}
\usepackage{enumitem}
\setlist{nosep}
\newtheorem{theorem}{Theorem}[section]
\newtheorem{lemma}[theorem]{Lemma}
\newtheorem{corollary}[theorem]{Corollary}
\newtheorem{proposition}[theorem]{Proposition}
\theoremstyle{remark}\newtheorem{remark}[theorem]{Remark}
\newcommand{\R}{\mathbb R}
\newcommand{\C}{\mathbb C}
\newcommand{\e}{\mathrm e}
\newcommand{\dd}{\,\mathrm d}
\newcommand{\ord}{\operatorname{ord}}
\newcommand{\supp}{\operatorname{supp}}
\newcommand{\GRH}{\mathrm{GRH}}
\newcommand{\abs}[1]{\lvert#1\rvert}
\newcommand{\N}{\mathbb N}
\title{Sharp support constants for positive filters of\newline Riemann and Dirichlet beta zeros}
\author{Research note for the ProveIt project}
\date{27 November 2026}
\begin{document}
\maketitle
\begin{abstract}
Let $H_\zeta(B)$ be the least half-width of a compactly supported probability measure whose Fourier transform vanishes at all nonzero ordinates of zeta zeros up to height $B$, with their multiplicities. Taking the vendored Lean theorem of RH as input, we show
\[H_\zeta(B)=\tfrac12\log B+O(\log\log B).\]
For the combined zero divisor of $\zeta$ and $L(s,\chi_{-4})$ the constant is $1$ instead of $1/2$. More generally it is $(r+s)/2$ for their fixed nonnegative integer powers $r,s$. These constants are sharp by Jensen's formula. The upper bound comes from a positive quadrature theorem for Fourier jets of a double-exponentially decaying measure. Its proof uses a separation-free extrapolation bound for exponential polynomials and a convex-hull argument, retaining exact zeros even at multiple or arbitrarily close ordinates. We also obtain positive filters with rational support locations and simultaneous uniform averaging of normalized prime-power errors over scales $X(\log X)^{\pm(1+o(1))}$ for the two residue classes modulo four. The new proofs are classical mathematical proofs, not Lean-checked proofs.
\end{abstract}

\section{Input, relation to the article, and statements}
The supplied repository contains a readable vendoring of the Lean development \cite{vendor,leanmain,leanstatements}. Its README pins upstream commit \texttt{72fead102f663d5c7a407964a13dc0cf9d977207} of \texttt{lsernau/grh-lean}, with Lean v4.35.0 and a separate Mathlib revision. The source \texttt{src/Main.lean} exports \texttt{riemannHypothesis}, \texttt{riemannZeta\_eq\_zero\_re}, \texttt{LFunction\_ne\_zero\_of\_}\allowbreak\texttt{half\_lt\_re}, \texttt{LFunction\_eq\_zero\_re}, and its primitive-character specialization. In particular the last statement about zeros in the open strip is for Dirichlet characters, not general automorphic or Artin $L$-functions. The README reports independent kernel checks; we inspected the statements and the reduction file, but did not rerun the full build or certify those reports independently. We use RH and Dirichlet GRH in precisely the scope of those statements. The positive quadrature argument below is itself unconditional.

The source's analytic reduction \cite{onesided} uses a ramp statistic with weight $(1/n-1/x)_+$ (interpreted as zero for $n>x$); a one-sided bound yields nonvanishing in the right half-plane. The merged article \cite{article,rec,zer} develops the converse spectral perspective: under Dirichlet GRH, normalized ramp errors have absolutely summable Fourier expansions on real zero ordinates. In particular, its \texttt{rec:thm:finitefilter}, \texttt{rec:thm:logfilter} construct positive filters of half-width $O_{\mathcal F}(\log B)$ for finite families and \texttt{rec:lem:jensensupport} gives the lower bound of order $\log B$. The article does not give a matching leading upper constant in these statements. We close that gap for two completed functions with positive theta kernels. Our proof is not an application of its separated-frequency coloring construction, and does not need local occupancy estimates for the upper bound.

\subsection*{Notation and main theorem}
Write $\beta(s)=L(s,\chi_{-4})$, where $\chi_{-4}(n)$ is $0$ on even integers, $1$ for $n\equiv1\pmod4$, and $-1$ for $n\equiv3\pmod4$. Let $N_j(B)$ count positive ordinates, with multiplicity, of nontrivial zeros of $j\in\{\zeta,\beta\}$ up to $B$. Under the input RH and Dirichlet GRH, the nontrivial zeros are $1/2+i\gamma$. At a shared ordinate define $m_{r,s}(\gamma)=r m_\zeta(\gamma)+s m_\beta(\gamma)$ for fixed integers $r,s\ge0$, $r+s>0$. Let $H_{r,s}(B)$ be the infimum of the half-widths of even probability measures $\mu$ on $[-H,H]$ with
\[
 \widehat\mu^{(k)}(\gamma)=0\qquad(0<\abs\gamma\le B,\ 0\le k<m_{r,s}(\gamma)),
 \qquad \widehat\mu(z)=\int \e^{izu}\dd\mu(u).
\]
Only nonzero multiplicities impose constraints. These derivatives are in the complex Fourier variable. Both completions will be shown nonzero at the center.

\begin{theorem}[Sharp leading support constant, using Dirichlet GRH]\label{thm:main}
For $B\to\infty$ and fixed $r,s$,
\begin{align*}
 H_{r,s}(B)&=\frac{r+s}{2}\log B+O_{r,s}(\log\log B),\\
 H_{r,s}(B)&\ge \frac{r+s}{2}\log B-\frac{r+s}{2}\log(2\pi)
 +\frac{s}{2}\log4-(r+s)-O_{r,s}(\log^2 B/B),\\
 H_{r,s}(B)&\le\frac{r+s}{2}(\log B+\log\log B)+O_{r,s}(1).
\end{align*}
The upper bound has a finite atomic witness. For $r=1,s=0$, one may take
\begin{equation}\label{eq:explicit}
 H=\tfrac12\log\bigl(256(2N_\zeta(B)+1)\bigr)
\end{equation}
with at most $N_\zeta(B)+1$ pairs of symmetric atoms. For one copy of either function the radii of the atoms can be chosen rational, with real nonnegative weights. No lower bound on differences of ordinates is used.
\end{theorem}
A ``pair'' may degenerate to one atom at zero. The lower bound concerns multiplicity-preserving cancellation, not simple cancellation of the distinct union. In particular it does not assume that the two functions have no common zeros. The upper bound also applies unconditionally to the real zeros of the two completed functions in the Fourier variable; the stated sharp asymptotics for the full ordinate multiset use the input theorem.

\section{Extrapolation without a separation hypothesis}
The following elementary estimate is related to Tur\'an's inequality for exponential polynomials \cite{nazarov}, but we include a proof for repeated exponents and the $L^1$ form we need.

\begin{lemma}[Uniform extrapolation]\label{lem:extra}
Let $f(t)=\sum_{j=1}^J P_j(t)\e^{i\lambda_j t}$, with real $\lambda_j$, $\deg P_j<m_j$, and $n=\sum_jm_j\ge1$. For $L>0$ and $t\ge L$,
\[
 |f(t)|\le \frac{1}{L}\left(A\frac{t}{L}\right)^n\int_0^L |f(v)|\dd v,\qquad A=64\e.
\]
There is no dependence on the sizes or mutual distances of the $\lambda_j$.
\end{lemma}
\begin{proof}
For $x,h\in\R$, put $u_k=f(x+kh)$, $k\ge0$. Let $r_j=\e^{i\lambda_jh}$ and $Q(z)=\prod_j(1-r_j z)^{m_j}=\sum_{\ell=0}^n q_\ell z^\ell$. Each summand of $u_k$ is a polynomial in $k$ of degree $<m_j$ times $r_j^k$, so the sequence obeys the recurrence with characteristic polynomial $\prod_j(E-r_j)^{m_j}$, where $E$ is forward shift. This works even when sampled roots $r_j$ coincide. Therefore, in formal power series,
\[
 Q(z)\sum_{k\ge0}u_kz^k=P(z),\qquad \deg P<n.
\]
If $M=\max_{0\le j<n}|u_j|$, then $|q_\ell|\le\binom n\ell$ and every coefficient of $P$ has modulus at most $2^n M$. Expanding each factor of $1/Q$ as a geometric series, with repetitions, gives $|[z^k]Q^{-1}|\le\binom{n+k-1}{n-1}$. Consequently
\[
 |u_k|\le n2^n\binom{n+k-1}{n-1}M
 \le \bigl(4\e(1+k/n)\bigr)^n M.
\]
The second inequality uses $n\le2^n$ and $\binom{n+k}n\le(\e(n+k)/n)^n$.

Now choose $k=\lceil2nt/L\rceil$. For each $x\in[0,L/2]$ set $h=(t-x)/k$. Then $x+kh=t$, and the first $n$ points $v_j(x)=x+j(t-x)/k$ are in $[0,L]$. Also $v_j'(x)=1-j/k\ge1/2$. Thus
\[
 \int_0^{L/2}\max_{j<n}|f(v_j(x))|\dd x
 \le 2n\int_0^L |f(v)|\dd v.
\]
Integrating the sequence bound over $x$ and dividing by $L/2$, we get $4n/L$ times $(4\e(1+k/n))^n$ times the integral. Since $1+k/n\le4t/L$ and $4n\le4^n$, the claimed estimate follows.
\end{proof}

\section{A compact positive Fourier-jet quadrature theorem}
\begin{theorem}[Double-exponential compactification]\label{thm:quad}
Let $w$ be an even probability density on $\R$ with, for $t\ge0$,
\begin{equation}\label{eq:env}
 c\e^{-a\e^{2t}}\le w(t)\le C\e^{-b\e^{2t}},\qquad a,b,c,C>0.
\end{equation}
Suppose $v_1,\ldots,v_d$ are real even exponential polynomials and all lie, together with the constant $1$, in a common exponential-polynomial space as in Lemma~\ref{lem:extra}, of order $n$. There is an even atomic probability measure $\nu$ supported on $[-H,H]$ with
\[
 H=\tfrac12\log(n+1)+O_w(1),\qquad
 \int v_j\dd\nu=\int v_j(t)w(t)\dd t\quad(1\le j\le d).
\]
At most $d+1$ symmetric pairs of atoms are needed, with rational radii. More precisely, if every nonzero scalar linear combination of $v_j-\int v_jw$ has order at most $n$ in a common space (not necessarily containing $1$), the same bound holds with that $n$.
\end{theorem}
\begin{proof}
We show that no nonzero real even exponential polynomial $f$ of order at most $n$, satisfying $\int_{\R} f w=0$, can be nonnegative on $[0,H]$. Let $H=\tfrac12\log(D(n+1))$, where $D$ is a constant to be fixed depending only on the envelope, and put $L=H/2\ge1$. If $f\ge0$ on $[0,H]$ and $f\not\equiv0$, then $I=\int_0^L f(t)\dd t>0$ by analyticity. Hence the positive part of the integral over $[0,H]$ is at least
\[
 \int_0^H f(t)w(t)\dd t\ge cI\e^{-a\e^H}.
\]
By Lemma~\ref{lem:extra}, the absolute tail is bounded by
\begin{align*}
 \int_H^\infty |f(t)|w(t)\dd t
 &\le \frac{CI}{L}\int_H^\infty (At/L)^n\e^{-b\e^{2t}}\dd t\\
 &\le \frac{CI}{L}\frac{\exp(n\log(2A)-b\e^{2H})}{2b\e^{2H}-n}.
\end{align*}
Indeed for $t=H+y$, $y\ge0$, we have $\log(t/H)\le y$ since $H\ge1$, and $\e^{2t}\ge\e^{2H}(1+2y)$. Choose $D\ge\e^4$ so large that
\[
 a\sqrt D\le bD/2,\quad bD\ge 1,\quad
 bD/2>\log(2A)+\max(0,\log(C/c))+2.
\]
Then $a\e^H\le bD(n+1)/2$ and $2b\e^{2H}-n\ge1$. Substitution shows the tail is strictly less than $cI\e^{-a\e^H}$, contradicting the zero integral by evenness. Replacing $f$ by $-f$ proves that every nonzero centered combination takes both signs on $[0,H]$.

Let $V(t)=(v_j(t)-\int v_jw)_{j=1}^d$ and $W$ be the real linear span of $\{V(t):0\le t\le H\}$. If $W=\{0\}$ the conclusion is immediate. Otherwise for every nonzero linear functional on $W$, the corresponding analytic $f(t)$ is not identically zero and takes both signs. Compact convex separation implies that $0$ is in the relative interior (in $W$) of the convex hull of $V([0,H])$. Here is also a justification of the rational-node assertion. On the unit sphere in $W$ the continuous function $a\mapsto\max_{0\le t\le H} a\cdot V(t)$ has a strictly positive minimum. By uniform continuity, a finite rational mesh in $[0,H]$ sufficiently fine has a positive maximum in every such direction. Its hull therefore contains $0$. Carath\'eodory's theorem expresses $0$ as a convex combination of at most $\dim W+1\le d+1$ mesh values. The symmetric probability measure placing half each weight at the two signed radii has the asserted moments. All inequalities concern exact functions; no numerical root separation enters.
\end{proof}

\begin{corollary}[Exact Fourier jets]\label{cor:jets}
For any finite list of real positive frequencies $\lambda$ and positive jet lengths $m(\lambda)$, the Fourier transform $F$ of $w$ can be matched, through order $m(\lambda)-1$, at all $\pm\lambda$ by an even compact probability measure of half-width
\[
 \tfrac12\log(2\sum_\lambda m(\lambda)+2)+O_w(1).
\]
If the listed jets of $F$ are all zero, one can use $2\sum_\lambda m(\lambda)$ as $n$ in Theorem~\ref{thm:quad}.
\end{corollary}
\begin{proof}
Use $v_{\lambda,k}(t)=t^k\cos(\lambda t+k\pi/2)$ for $0\le k<m(\lambda)$, which are real even functions. Differentiation under the integral is permitted by the envelope. These functions lie in the space generated by $t^k\e^{\pm i\lambda t}$; adding the constant increases its order by at most one. Evenness gives the negative-frequency constraints. When the target is zero there is no added constant.
\end{proof}

\begin{proposition}[Worst-case optimality of the general quadrature coefficient]\label{prop:worst}
For a fixed $w$ satisfying \eqref{eq:env}, the leading constant $1/2$ in the bound $H\le\tfrac12\log n+O_w(1)$ cannot be reduced for general even exponential-polynomial spaces of order $n$.
\end{proposition}
\begin{proof}
Take the polynomial space of degree at most $2m$, of order $n=2m+1$ (exponent $0$ with multiplicity $n$). Matching its even moment of degree $2m$ by a measure supported on $[-H,H]$ requires
$H\ge(\int |t|^{2m} w(t)\dd t)^{1/(2m)}$. For large $m$ take $R=(\log m-2\log\log m)/2>0$. Integration over $[R,R+1]$ and the lower envelope give
\[
 H\ge R\left(c\exp(-a\e^2 m/(\log m)^2)\right)^{1/(2m)}
 =(\tfrac12-o(1))\log m.
\]
Evenness of $w$ and of the matching measure makes odd moments automatic. This yields the stated leading coefficient for $n=2m+1$.
\end{proof}

\section{Positive kernels of the two completed functions}
The Fourier normalization is important: the variable $t$ in $\Xi(t)=\xi(1/2+it)$ is the actual zeta ordinate (not half an ordinate).

\subsection{The Riemann kernel}
Let $\psi(x)=\sum_{k\ge1}\e^{-\pi k^2x}$ and $g(u)=\e^{u/2}\psi(\e^{2u})$. The standard Mellin and Jacobi identities give \cite{titchmarsh,edwards}
\[
 \xi(s)=\tfrac12s(s-1)\pi^{-s/2}\Gamma(s/2)\zeta(s)
 =\tfrac12+\tfrac{s(s-1)}2\int_1^\infty\psi(x)(x^{s/2}+x^{(1-s)/2})\frac{\dd x}{x}.
\]
For clarity, splitting the defining Mellin integral at $1$ and using $1+2\psi(x)=x^{-1/2}(1+2\psi(1/x))$ yields this identity first for $\Re s>1$ and then everywhere by continuation. The same identity implies $g(-u)=g(u)+(\e^{u/2}-\e^{-u/2})/2$ and thus $g'(0)=-1/4$. Define $K_\zeta(u)=g''(u)-g(u)/4$ for $u\ge0$ and extend evenly. Twice integrating by parts, with the rapid decay at infinity, gives
\begin{align*}
 \Xi(t)&=2\int_0^\infty K_\zeta(u)\cos(tu)\dd u
 =\int_{\R}K_\zeta(u)\e^{itu}\dd u,\\
 K_\zeta(u)&=\sum_{k\ge1}\left(4\pi^2k^4\e^{9u/2}-6\pi k^2\e^{5u/2}\right)\e^{-\pi k^2\e^{2u}}.
\end{align*}
All termwise differentiations are locally uniformly convergent. The summands are positive on $u\ge0$ since $2\pi k^2\e^{2u}>3$. In particular $\Xi(0)>0$. With $x=\e^{2u}\ge1$,
\begin{equation}\label{eq:zetaenv}
 \e^{-\pi x}\le K_\zeta(u)\le 8\pi^2\e^{-x/2}.
\end{equation}
For the lower bound retain $k=1$, whose polynomial-exponential prefactor is at least $2\pi(2\pi-3)>1$. For the upper bound discard the negative terms; $\sum_{k\ge1} k^4\e^{-\pi(k^2-1)x}<2$ because $4\log k\le2(k^2-1)$ and $\pi-2>1$, with the terms for $k\ge2$ bounded by $\e^{-k}$. Finally $x^{9/4}\e^{-(\pi-1/2)x}\le1$ by $\log x\le x$ and $\pi-1/2>9/4$. Normalizing $w_\zeta=K_\zeta/\Xi(0)$ gives a probability density satisfying \eqref{eq:env}, with $a=\pi,b=1/2,C/c=8\pi^2$.

In Theorem~\ref{thm:quad} one may use $D=256$: $a\sqrt D=16\pi<64=bD/2$, $D>\e^4$, and $64>\log(128\e)+\log(8\pi^2)+2$. This proves the explicit constant in \eqref{eq:explicit} for $n=2N_\zeta(B)$, if the cutoff includes at least one zero. Otherwise the unit atom at zero suffices. The same explicit bound for a chosen multiset of real $\Xi$-zeros is unconditional.

\subsection{The Dirichlet beta kernel}
Define
\[
 \mathcal L_\beta(s)=(4/\pi)^{(s+1)/2}\Gamma((s+1)/2)\beta(s),\qquad
 \theta_\beta(x)=\sum_{k\ge0}(-1)^k(2k+1)\e^{-\pi(2k+1)^2x/4}.
\]
The odd primitive-character theta transformation (Poisson summation, Gauss sum $2i$, root number $1$) gives $\theta_\beta(x)=x^{-3/2}\theta_\beta(1/x)$; see \cite{kedlaya}. For $\Re s>1$, termwise Mellin integration yields $\mathcal L_\beta(s)=\int_0^\infty \theta_\beta(x)x^{(s+1)/2}\dd x/x$. The transformation provides rapid decay at zero, so this identity continues to all $s$. With $x=\e^{2u}$ it becomes
\[
 \mathcal L_\beta(1/2+it)=\int_{\R}K_\beta(u)\e^{itu}\dd u,
 \qquad K_\beta(u)=2\e^{3u/2}\theta_\beta(\e^{2u}).
\]
The theta transformation makes $K_\beta$ even. For $x\ge1$, the positive terms $(2k+1)\e^{-\pi(2k+1)^2x/4}$ decrease strictly to zero: the real function $y\e^{-\pi y^2x/4}$ decreases for $y\ge1$. The second-to-first term ratio is $3\e^{-2\pi x}<1/2$. Alternating-series bounds therefore give
\[
 \tfrac12\e^{-\pi x/4}\le\theta_\beta(x)\le\e^{-\pi x/4}.
\]
Consequently $K_\beta$ is strictly positive, $\mathcal L_\beta(1/2)>0$, and
\[
 \e^{-\pi x/4}\le K_\beta(u)\le C_\beta\e^{-\pi x/8}\quad(u\ge0),
 \qquad C_\beta=2\sup_{x\ge1}x^{3/4}\e^{-\pi x/8}<\infty.
\]
Normalize by $\mathcal L_\beta(1/2)$ to obtain $w_\beta$ satisfying \eqref{eq:env}. This is an instance of the positive-theta-kernel condition, not a claim that every Dirichlet completion has such a positive kernel.

\section{Proof of the sharp support constants}
\begin{proof}[Proof of Theorem~\ref{thm:main}]
By the supplied RH and Dirichlet GRH, the nontrivial zeros of both completions are real zeros in their Fourier variable $t$. The gamma factors introduce no extra zeros of the completed functions in the open strip, and trivial zeros have been canceled in the completions. We have also proved both central values positive. For each $j=\zeta,\beta$ apply Corollary~\ref{cor:jets} to all positive zeros up to $B$ with multiplicity. Obtain an even atomic measure $\mu_j$ of half-width
\[
 h_j\le\tfrac12\log(2N_j(B)+1)+O_j(1)
\]
with at least the required vanishing orders. Use the convolution $\mu_\zeta^{*r}*\mu_\beta^{*s}$ (power zero means unit atom at zero). Its transform is the product, its vanishing orders add, and its support half-width is at most $rh_\zeta+sh_\beta$. The fixed-conductor Riemann--von Mangoldt formulas \cite{titchmarsh,davenport} are
\[
 N_j(T)=\frac{T}{2\pi}\log\frac{q_jT}{2\pi}-\frac{T}{2\pi}+O_j(\log(T+2)),
 \qquad q_\zeta=1,\quad q_\beta=4.
\]
Inserting these yields the stated upper bound. Finiteness and rational radii persist under convolution. The explicit zeta bound was proved above.

For the lower bound, let $F(z)=\widehat\mu(z)$ for any admissible probability measure supported on $[-H,H]$. It is entire, $F(0)=1$, and $|F(Re^{i\theta})|\le\exp(HR|\sin\theta|)$. Jensen's formula bounds the sum of $\log(R/|z_k|)$ over zeros in the disk (with multiplicity) by $2HR/\pi$. Keeping just the prescribed symmetric zeros and letting $R\downarrow B$ gives
\begin{equation}\label{eq:jensen}
 H\ge\frac\pi B\int_0^B\frac{rN_\zeta(t)+sN_\beta(t)}{t}\dd t.
\end{equation}
A zero exactly at $B$ contributes zero to this integral, so endpoint conventions do not matter. The finite integral near zero is harmless because the central values are nonzero and the zero sets are discrete. For each $j$, integrate the displayed zero-count main term on $[2,B]$ and absorb the lower range in $O(1)$. This gives
\[
 \frac\pi B\int_0^B\frac{N_j(t)}t\dd t
 =\tfrac12\log B+\tfrac12\log(q_j/(2\pi))-1+O_j(\log^2 B/B).
\]
Together with the upper bound this proves all assertions.
\end{proof}
\begin{remark}
In fact the infimum defining $H_{r,s}(B)$ is attained: take feasible half-widths tending to the infimum and use weak compactness of probability measures in one common compact interval. The finitely many Fourier-jet conditions are continuous against weak convergence, and the limiting measure is supported on the limiting interval. Symmetry is preserved. This does not assert uniqueness or exact optimal finite-$B$ width.
\end{remark}

\section{Uniform averaging consequences for primes}
\subsection{Adding tail attenuation at constant support cost}
Let $U_B$ be the uniform probability measure on $[-1/B,1/B]$. Its transform is $\sin(z/B)/(z/B)$, interpreted as $1$ at zero. For real $|z|\ge B$,
$|\widehat U_B(z)|\le\sin1=:\sigma<1$: on $[1,\pi]$ the sinc function is positive and decreasing, and beyond $\pi$ its modulus is at most $1/\pi<\sin1$. Set $c_0=-\log\sigma$. Convolve a filter from Theorem~\ref{thm:main} with $k$ copies of $U_B$. The exact vanishing orders are retained, the support half-width increases by at most $k/B$, and the transform on $|z|\ge B$ has modulus at most $\e^{-c_0 k}$.

\begin{corollary}[Sharp leading width with exponentially small leakage]\label{cor:leak}
For fixed $r,s$ and any fixed $\alpha>0$ there are even positive probability measures with the exact zero-jet constraints of Theorem~\ref{thm:main}, half-width
\[
 H\le\frac{r+s}{2}(\log B+\log\log B)+O_{r,s,\alpha}(1),
\]
and $\sup_{|\omega|\ge B}|\widehat\mu(\omega)|\le\e^{-\alpha B}$. The leading support constant remains optimal for this stronger class of filters.
\end{corollary}
\begin{proof}
Take $k=\lceil\alpha B/c_0\rceil$, so $k/B=O_\alpha(1)$. The lower bound follows because exact low zero cancellation is still required.
\end{proof}
For every absolutely convergent Fourier series $V(t)=b+\sum_{\gamma\ne0}a_\gamma\e^{i\gamma t}$ with frequencies drawn from the selected divisor, the uniform-in-center averaging error is at most $\e^{-\alpha B}\sum_{\gamma\ne0}\|a_\gamma\|$, in any finite-dimensional normed coefficient space. This follows by termwise integration. The result also retains exact cancellation of $t^j\e^{i\gamma t}$ for $j<m_{r,s}(\gamma)$ at low frequencies, by differentiating the transform and expanding $(t+u)^j$.

\subsection{Two prime races modulo four}
For $a=1,3$ set
\[
 S_a(x)=\sum_{\substack{n\le x\\ n\equiv a\ (4)}}\Lambda(n)(1/n-1/x),\qquad R_a(x)=S_a(x)-\tfrac12\log x-c_a,
\]
where $\Lambda$ is the von Mangoldt function and $c_a$ is the constant term of this ramp sum. The ramp explicit formula in \cite[\texttt{zer:lem:ramp}, \texttt{zer:eq:Zexp}]{zer} (or Mellin inversion of $-L'/L$ times $x^{s-1}/(s(s-1))$) gives under the input Dirichlet GRH
\begin{equation}\label{eq:ramp}
 \e^{t/2}R_a(\e^t)=\sum_{\gamma\ne0}\frac{m_\zeta(\gamma)+\chi_{-4}(a)m_\beta(\gamma)}{2(1/4+\gamma^2)}\e^{i\gamma t}+O((1+t)\e^{-t/2}),\quad t\ge\log2.
\end{equation}
The principal character modulo four is induced from zeta, and its omitted Euler factor at $2$ changes only $c_a$ and the indicated error. There is no central term because the positive kernels show both central values nonzero. The coefficients are absolutely summable by the fixed-conductor zero counts. To recall the derivation: on a line $c>1$ the integral equals $S_\chi(x)$ by the ramp Mellin kernel. Moving to $\Re s=-1/2$ crosses the pole at $1$ for the principal inducer, the nontrivial zeros giving $m_\gamma x^{-1/2+i\gamma}/(1/4+\gamma^2)$, and at most a double pole at $0$ with residue $O((1+\log x)/x)$. The functional equation bounds the new vertical integral by $O(x^{-3/2})$ and standard zero-avoiding heights make the horizontal integrals vanish. Removing the powers of $2$ changes the constant and $O((1+\log x)/x)$; orthogonality of the two characters gives \eqref{eq:ramp}.

\begin{theorem}[A common near-optimal multiplicative window, using Dirichlet GRH]\label{thm:prime}
For large $X$ put $T=\log X$. There is a common even probability measure $\mu_X$ independent of the class $a$, supported on $[-H_X,H_X]$ with
\[
 H_X\le\log T+\log\log T+O(1),
\]
for which
\[
 \max_{a=1,3}\left|\int \e^{u/2} R_a(X\e^u)\dd\mu_X(u)\right|
 \ll \frac{T^{3/2}(\log T)^{1/2}}{X}.
\]
Thus all scales lie between $X/(C\log X\log\log X)$ and $CX\log X\log\log X$ for a fixed $C>0$. The same averaging rule can be used at every center $Y$ with $\log Y$ comparable to $T$, with the bound in the proof. For the zeta ramp alone the half-width can be taken $\tfrac12(\log T+\log\log T)+O(1)$.
\end{theorem}
\begin{proof}
Take $r=s=1$, $B=T$ and $\alpha=1$ in Corollary~\ref{cor:leak}. The spectral sum of \eqref{eq:ramp}, averaged at any center $v$ with $v-H_X\ge\log2$, is bounded by $C_1\e^{-T}$. Its remainder is bounded by $C_2(1+v+H_X)\e^{-(v-H_X)/2}$. At $v=\log X=T$, multiplication by $X^{-1/2}$ gives $O(X^{-3/2}+(1+T+H_X)\e^{H_X/2}/X)$, as required. The bound at $v=\log Y$ is the same before replacing $v$ by $T$; no center-dependent choice of filters was needed. With zeta alone use $r=1,s=0$ and the corresponding single-function ramp formula.
\end{proof}
This is a statement about smoothing prime-power errors, not an improvement of the unsmoothed pointwise prime number theorem. The sharpness here is for the underlying exact multiplicity-preserving spectral filter; the smallest window for a prime-error estimate alone is a different optimization problem.

\section{Novelty, limitations, and further questions}
\paragraph{What is added.} The article's \texttt{rec:thm:logfilter} supplies a logarithmic order via coloring and convolution and its Jensen lemma supplies a leading lower constant. Theorem~\ref{thm:main} supplies the matching leading upper constant for the zeta and beta divisors, with explicit zeta bound, exact jets, no spacing assumption, rational support locations, and an optimal leading coefficient even with exponentially small leakage. The underlying positive quadrature theorem applies to arbitrary real Fourier jets (including polynomial moments), not only the zero ordinates. The classical ingredients---theta kernels, Jensen's formula, zero counting, exponential-polynomial extrapolation and convex separation---are not claimed new. A targeted literature search for their combination in this exact support problem did not establish priority; this note claims a proposed original extension relative to the inspected article, not a certified first discovery in the literature.

\paragraph{What is not claimed.} We have not independently rebuilt the 818-file vendored package, checked every theorem in the 723-page merged article, or formalized the new argument. No result for arbitrary Dirichlet characters with possibly signed theta kernels follows automatically, nor for nonabelian $L$-functions. No claim about simplicity, independence or mutual disjointness of zeros is made. The finite witnesses are existence proofs; the accompanying linear-programming experiments use approximate zero ordinates and floating point, so they are not exact certificates. The main proof needs no experimental input.

\paragraph{Research questions.} (1) Can the $\frac12\log\log B$ excess in the upper bound be removed, or is there a genuine second-order term? (2) Which other completed Dirichlet functions have positive kernels, or admit multiplication by an auxiliary Fourier transform of small exponential type to obtain positivity? (3) Can one bound the rational mesh size and weights explicitly using zero enclosures and multiplicities, without a zero-spacing conjecture? (4) Can the optimal leading constant for cancellation of distinct ordinates alone be determined, including the effect of multiplicities and shared zeros? (5) Does the optimal width for approximate cancellation alone have a universal leading coefficient under a natural robustness specification? (6) Can the compactification lemma be formalized using Mathlib's convex hull, finite-dimensional separation and formal power series, then linked to the pinned Lean GRH theorem? The dependency chain is recurrence $\to$ extrapolation $\to$ envelope tail $\to$ convex hull $\to$ positive quadrature $\to$ Jensen and zero counts.

\begin{thebibliography}{99}
\bibitem{vendor} ProveIt, vendored \texttt{grh-lean} README, retrieved 27 Nov. 2026. \href{https://github.com/VladimirReshetnikov/ProveIt/blob/main/Analysis/GeneralizedRiemannHypothesis/vendor/grh-lean/README.md}{Repository file}. Upstream commit \texttt{72fead102f663d5c7a407964a13dc0cf9d977207}.
\bibitem{leanmain} \texttt{grh-lean/grh/src/Main.lean}, \href{https://github.com/VladimirReshetnikov/ProveIt/blob/main/Analysis/GeneralizedRiemannHypothesis/vendor/grh-lean/grh/src/Main.lean}{vendored theorem exports}.
\bibitem{leanstatements} \texttt{grh-lean/grh/src/Statements.lean}, \href{https://github.com/VladimirReshetnikov/ProveIt/blob/main/Analysis/GeneralizedRiemannHypothesis/vendor/grh-lean/grh/src/Statements.lean}{formal propositions}.
\bibitem{onesided} \texttt{grh-lean/grh/src/New/Reduction/OneSidedBound.lean}, \href{https://github.com/VladimirReshetnikov/ProveIt/blob/main/Analysis/GeneralizedRiemannHypothesis/vendor/grh-lean/grh/src/New/Reduction/OneSidedBound.lean}{analytic reduction}.
\bibitem{article} ProveIt Contributors, \emph{Beyond the Riemann hypothesis for Dirichlet $L$-functions: consequences, transfers and limits}, 26 Nov. 2026. \href{https://github.com/VladimirReshetnikov/ProveIt/tree/main/Analysis/GeneralizedRiemannHypothesis/Article}{Article collection}.
\bibitem{rec} Ibid., \href{https://github.com/VladimirReshetnikov/ProveIt/blob/main/Analysis/GeneralizedRiemannHypothesis/Article/sections/rec.tex}{\texttt{sections/rec.tex}}, particularly \texttt{rec:thm:finitefilter}, \texttt{rec:thm:logfilter}, \texttt{rec:lem:jensensupport}, \texttt{rec:thm:uniformavg}, \texttt{rec:thm:polylogavg}.
\bibitem{zer} Ibid., \href{https://github.com/VladimirReshetnikov/ProveIt/blob/main/Analysis/GeneralizedRiemannHypothesis/Article/sections/zer.tex}{\texttt{sections/zer.tex}}, ramp explicit formula and spectral averaging.
\bibitem{titchmarsh} E. C. Titchmarsh, revised by D. R. Heath-Brown, \emph{The Theory of the Riemann Zeta-function}, 2nd ed., Oxford University Press, 1986.
\bibitem{edwards} H. M. Edwards, \emph{Riemann's Zeta Function}, Academic Press, 1974; Dover reprint, 2001.
\bibitem{davenport} H. Davenport, revised by H. L. Montgomery, \emph{Multiplicative Number Theory}, 3rd ed., Springer, 2000.
\bibitem{kedlaya} K. S. Kedlaya, \emph{Notes on analytic number theory}, \href{https://kskedlaya.org/ant/chap-funceq2.html}{Functional equations for Dirichlet $L$-functions}, section 6.2 (odd characters).
\bibitem{nazarov} F. L. Nazarov, Local estimates for exponential polynomials and their applications to inequalities of the uncertainty principle type, \emph{St. Petersburg Math. J.} 5 (1994), 663--717. (Background; our extrapolation lemma is proved directly.)
\end{thebibliography}
\end{document}
